\stfhSubIndex{I080}{lian^luo}{联络（connection）}{dian^fan^lian^luo}{典范联络（canonical connection）}
\stfhIndex{I178}{xiao^lin^zhao^qi}{小林昭七（Kobayashi）}
本章将在殆凯勒流形 $(M,\omega,J)$ 的背景下，对一般伪全纯曲线作几何计算。我们先在实情形中，推导可能带有拉格朗日边界的伪全纯曲线之调和能量密度的一些一般事实。这些几何计算的一个主要推论，是命题 7.4.5 给出的、具有拉格朗日边界条件的 $J$-全纯曲线之先验 $W^{2,2}$ 估计。

在本节关于能量密度的计算中，我们将利用由夏尔·埃雷斯曼（Charles Ehresmann）与波莱特·利贝尔曼（Paulette Libermann） (EL51) 引入、并由小林昭七（Shoshichi Kobayashi；日文：小林昭七） (Ko03) 研究的殆凯勒流形全纯切丛之\emph{典范联络}概念。（关于殆复流形上各种自然联络的广泛研究，另见 (Ga97)。）我们将看到，在全纯切丛上使用这个典范联络，会极大简化涉及 $J$-全纯映射能量密度之导数的计算。虽然这种简化对椭圆估计并非必不可少，但我们更愿意采用这种最优形式的计算，并希望它将来会有用。为了最大限度地发挥典范联络的优势，在局部计算中，我们将对殆凯勒流形 $(M,\omega,J)$ 和黎曼曲面 $(\Sigma,j)$ 采用复化记号。

\section{殆凯勒流形上的自然联络}
\label{sec:natural-connection-almost-kahler}

本节先回顾由埃雷斯曼与利贝尔曼引入，并由保罗·戈迪雄（Paul Gauduchon） (Ga97) 和小林昭七 (Ko03) 研究的典范联络概念。

我们先描述殆埃尔米特流形 $(M,J,g)$ 上典范联络的基本性质，再描述殆凯勒流形上的情形。随后，当 $(M,J,g)$ 是殆凯勒的，即基本 $2$-形式 $\Phi=g(\cdot,J\cdot)$ 闭时，我们从 $(M,g)$ 的列维--奇维塔联络出发，给出典范联络的构造（见定理 7.1.14）。特别地，后一构造适用于三元组 $(M,\omega,J)$，其中 $(M,\omega)$ 是辛流形，且 $J$ 与 $\omega$ 相容。

\subsection{典范联络}
\label{subsec:canonical-connection}

设 $(M,J)$ 是殆复流形，即 $J$ 是恒等映射之上的丛同构
\[
 J:TM\longrightarrow TM,
 \qquad J^2=-id.
\]

\begin{definition}\label{def:hermitian-metric-almost-complex}
若 $(M,J)$ 上的度量 $g$ 满足
\[
 g(Ju,Jv)=g(u,v),\qquad u,v\in T_xM,\quad x\in M,
\]
\stfhIndex{I054}{dai^ai^er^mi^te^liu^xing}{殆埃尔米特流形（almost-Hermitian manifold）}
则称 $g$ 是埃尔米特的。我们称三元组 $(M,J,g)$ 为\emph{殆埃尔米特流形}。
\end{definition}

对任意给定的殆埃尔米特流形 $(M,J,g)$，双线性形式
\[
 \Phi:=g(J\cdot,\cdot)
\]
在 (KN96) 中称为\emph{基本二形式}；它是非退化的。

\begin{definition}\label{def:almost-kahler-manifold}
\stfhIndex{I055}{dai^kai^lei^liu^xing}{殆凯勒流形（almost-Kähler manifold）}
若上述二形式 $\Phi$ 是闭的，则称殆埃尔米特流形 $(M,J,g)$ 为\emph{殆凯勒流形}。
\end{definition}

\begin{definition}\label{def:almost-hermitian-connection}
\stfhSubIndex{I081}{lian^luo}{联络（connection）}{ai^er^mi^te^lian^luo}{埃尔米特联络（Hermitian connection）}
（殆）埃尔米特联络 $\nabla$ 是满足
\[
 \nabla g=0=\nabla J
\]
的仿射联络。
\end{definition}

\begin{exercise}\label{ex:hermitian-connection-exists}
证明这样的联络总是存在。
\end{exercise}

一般而言，殆埃尔米特联络 $\nabla$ 的挠率 $T=T_\nabla$ 不为零，即使 $J$ 可积时也是如此。依定义，$T$ 的表达式为
\[
 T(X,Y)=\nabla_XY-\nabla_YX-[X,Y].
\]
下面是复几何中陈联络的殆复版本。

\begin{theorem}[{(Ga97), (Ko03)}]\label{thm:unique-hermitian-canonical-connection}
在任意殆埃尔米特流形 $(M,J,g)$ 上，存在唯一的 $TM$ 上埃尔米特联络 $\nabla$，对所有 $X\in TM$ 满足
\begin{equation}
 T(X,JX)=0.
 \label{eq:canonical-torsion-property}
\end{equation}
\end{theorem}

在 $J$ 可积的复几何 (Cher67) 中，满足\eqref{eq:canonical-torsion-property}的埃尔米特联络称为陈联络。

\begin{definition}\label{def:canonical-connection-torsion}
殆埃尔米特联络的一个\emph{典范联络}，定义为具有挠率性质\eqref{eq:canonical-torsion-property}的联络。
\end{definition}

接着，我们希望推导双线性形式 $\Phi=g(\cdot,J\cdot)$ 的殆凯勒性质 $d\Phi=0$ 如何影响典范联络的挠率性质。为此，最好考虑 $\nabla$ 的复化。考虑复化切丛
\[
 T_{\mathbb C}M=TM\otimes\mathbb C
 =T^{(1,0)}M\oplus T^{(0,1)}M,
\]
其中 $T^{(1,0)}M$（相应地，$T^{(0,1)}M$）是 $J$ 对应于特征值 $i=\sqrt{-1}$（相应地，特征值 $-i$）的特征空间。度量 $g$ 复线性地延拓为 $T_{\mathbb C}M$ 上的非退化复双线性形式，仍记为 $g$；关于它，$T^{(1,0)}M$ 与 $T^{(0,1)}M$ 都成为迷向的，即形式 $\Phi$ 分别限制在 $T^{(1,0)}M$ 和 $T^{(0,1)}M$ 上时消失。于是，复向量空间 $T^{(1,0)}M$ 上相伴的埃尔米特内积定义为
\begin{equation}
 \langle u,v\rangle=g(u,\overline v)
 \label{eq:hermitian-inner-product-complexified}
\end{equation}
这里，$u,v\in T_x^{(1,0)}M$ 是两个复向量。

现在，相对于给定标架来表示这样的 $\nabla$ 之复化以及相伴的 $T$。在 $TM$ 上选取 $M$ 的一个正交归一标架
\begin{equation}
 \{e_1,\ldots,e_n,f_1,\ldots,f_n\}
 \label{eq:real-orthonormal-frame}
\end{equation}
并使 $f_j=Je_j$。于是向量
\[
 u_j=\frac1{\sqrt2}(e_j-if_j),\qquad j=1,\ldots,n,
\]
构成 $T^{(1,0)}M$ 的一个酉标架，而向量
\stfhIndex{I296}{you^biao^jia}{酉标架（unitary frame）}
\[
 \overline u_j=\frac1{\sqrt2}(e_j+if_j),\qquad j=1,\ldots,n,
\]
构成 $T^{(0,1)}M$ 的一个酉标架。设
\begin{equation}
 \{\alpha^1,\ldots,\alpha^n,\beta^1,\ldots,\beta^n\}
 \label{eq:dual-real-frame}
\end{equation}
是\eqref{eq:real-orthonormal-frame}的对偶标架。复值 $1$-形式
\[
 \theta^j=\frac1{\sqrt2}(\alpha^j+i\beta^j),\qquad j=1,\ldots,n,
\]
构成 $\Lambda^{(0,1)}(M)=(T^{(1,0)}M)^*$ 的一个酉标架，它与酉标架 $\{u_1,\ldots,u_n\}$ 对偶。

$TM$ 上任意殆埃尔米特联络 $\nabla$ 都诱导复向量丛 $\Lambda^{(1,0)}M$ 上的一个联络；它由局部定义的矩阵值复 $1$-形式
$\omega=(\omega_j^i)_{i,j=1,2,\ldots,n}$ 定义，并且是斜埃尔米特的，即
$\overline{\omega_j^i}=-\omega_i^j$。相应的挠率记为 $\Theta$，它是取值于 $T^{(1,0)}M$ 的斜埃尔米特二形式，并具有表达式
\begin{equation}
 \Theta=\Pi'T_{\mathbb C},
 \label{eq:torsion-projection-complexification}
\end{equation}
其中 $T_{\mathbb C}=T\otimes\mathbb C$ 是 $T$ 的复线性延拓，而 $\Pi'$ 是到 $T^{(1,0)}M$ 的投影。

\stfhIndex{I126}{di^yi^jie^gou^fang^cheng}{第一结构方程（first structure equation）}
这个酉联络的第一结构方程为
\begin{equation}
 d\theta^i=-\sum_j\omega_j^i\wedge\theta^j+\Theta^i,
 \label{eq:first-structure-unitary-connection}
\end{equation}
用矩阵方程简写就是 $d\theta=-\omega\wedge\theta+\Theta$。
第二结构方程为
\stfhIndex{I244}{di^er^jie^gou^fang^cheng}{第二结构方程（second structure equation）}
\begin{equation}
 d\omega_j^i=-\sum_k\omega_k^i\wedge\omega_j^k+\Omega_j^i,
 \label{eq:second-structure-unitary-connection}
\end{equation}
用矩阵方程简写就是 $d\omega=-\omega\wedge\omega+\Omega$，其中 $\Omega$ 是曲率二形式。二者分解为
\[
 \Theta=\Theta^{(2,0)}+\Theta^{(1,1)}+\Theta^{(0,2)},
 \qquad
 \Omega=\Omega^{(2,0)}+\Omega^{(1,1)}+\Omega^{(0,2)}.
\]

\begin{theorem}\label{thm:torsion-zero-two-independent}
挠率形式 $\Theta$ 的 $(0,2)$-分量 $\Theta^{(0,2)}$ 与殆埃尔米特联络的选择无关。
\end{theorem}

\begin{exercise}\label{ex:torsion-zero-two-independent}
证明这个定理。
\end{exercise}

这个 $(0,2)$-形式 $\Theta^{(0,2)}$ 称为殆复结构 $J$ 的\emph{挠率}。$J$ 的可积性由 $\Theta^{(0,2)}=0$ 给出。另一方面，分量 $\Theta^{(1,1)}$ 和 $\Theta^{(2,0)}$ 依赖于联络的选择。下面是复化形式的典范联络定义，它等价于定义 7.1.6。

\begin{definition}\label{def:canonical-connection-complexified}
\stfhSubIndex{I080}{lian^luo}{联络（connection）}{dian^fan^lian^luo}{典范联络（canonical connection）}
设 $(M,J,g)$ 是殆埃尔米特流形。若 $TM$ 的联络 $\nabla$ 在复向量丛 $T^{(1,0)}M$ 上的挠率形式 $\Theta$ 满足 $\Theta^{(1,1)}=0$，则称 $\nabla$ 为典范联络。
\end{definition}

小林昭七证明了：若 $(M,J,g)$ 是殆凯勒的，则有下述结论。

\begin{theorem}[小林昭七，(Ko03)]\label{thm:almost-kahler-torsion-zero-two}
设 $(M,J,g)$ 是殆凯勒流形，$\nabla$ 是 $TM$ 的典范联络及其复化。则还成立 $\Theta^{(2,0)}=0$，因而 $\Theta$ 具有型 $(0,2)$，并满足
\begin{equation}
 \sum_i\Theta^i\wedge\overline\theta^i=0.
 \label{eq:torsion-cyclic-symmetry}
\end{equation}
\end{theorem}

等价地，$\Theta$ 具有型 $(0,2)$ 可由下述实描述给出。

\begin{proposition}\label{prop:canonical-torsion-real-identities}
设 $(M,J,g)$ 是殆凯勒流形，$\nabla$ 是其典范联络。则挠率张量 $T$ 满足
\[
 T(JY,Z)=T(Y,JZ),\qquad JT(JY,Z)=T(Y,Z)
\]
这里，$Y,Z$ 是任意向量场。
\end{proposition}

\begin{remark}\label{rem:canonical-connection-real-nonunique}
从实切丛 $TM$ 上的一般仿射联络出发时，会有许多不同的这种仿射联络，其复化都在全纯切丛 $T^{(1,0)}M$ 上给出满足 $\Theta^{(1,1)}=0$ 的联络。这样的仿射联络并不唯一。
\end{remark}

在下一小节中，我们将说明，除了上述典范埃尔米特联络外，列维--奇维塔联络也是这样一种选择。不过，这个列维--奇维塔联络并非 $J$-线性的。因此，我们倾向于使用上述定理给出的唯一埃尔米特联络。

\subsection{列维--奇维塔联络}
\label{subsec:levi-civita-connection}

回忆一下，把殆埃尔米特流形 $(M,J,g)$ 看作黎曼流形时，它在 $TM$ 上诱导列维--奇维塔联络，并可复线性地延拓到 $T_{\mathbb C}M$。

现在把这个列维--奇维塔联络复线性地延拓到 $T_{\mathbb C}M$，并用上述活动标架来表示它。$T_{\mathbb C}M$ 上的列维--奇维塔联络 $\nabla$ 相对于基
\[
 \{u_1,\ldots,u_n,\overline u_1,\ldots,\overline u_n\}
\]
局部表示为
\[
 \nabla u_j=\sum_i\lambda_j^i\otimes u_i
 +\sum_\ell\mu_j^{\overline\ell}\otimes\overline u_\ell
\]
及其复共轭。第一结构方程变成
\begin{equation}
 d\begin{pmatrix}\theta\\ \overline\theta\end{pmatrix}
 =-\begin{pmatrix}\lambda&\mu\\ \overline\mu&-\overline\lambda\end{pmatrix}
 \wedge\begin{pmatrix}\theta\\ \overline\theta\end{pmatrix}
 \wedge\begin{pmatrix}\theta\\ \overline\theta\end{pmatrix}.
 \label{eq:levi-civita-first-structure-source}
\end{equation}
由于 $\overline\lambda=-\lambda^t$，斜埃尔米特矩阵
\stfhSubIndex{I082}{lian^luo}{联络（connection）}{li^xi^nei^luo^wei^ci^lian^luo}{利希内罗维茨联络（Lichnerowicz connection）}
$\lambda=(\lambda_j^i)$ 在 $T^{(1,0)}M$ 上诱导一个联络，称为\emph{利希内罗维茨联络}。$\lambda$ 的第一结构方程为
\[
 d\theta^j=-\sum_i\lambda_i^j\wedge\theta^i
 -\sum_i\mu_{\overline i}^{j}\wedge\overline\theta^i.
\]
所以，回忆挠率形式的定义\eqref{eq:first-structure-unitary-connection}，$T^{(1,0)}M$ 上这个联络的挠率形式 $\Theta'$ 为
\begin{equation}
 (\Theta')^j=-\sum_i\mu_{\overline i}^{j}\wedge\overline\theta^i.
 \label{eq:lichnerowitz-torsion}
\end{equation}
我们写成
\[
 \mu_{\overline i}^{j}
 =\sum_{i,\ell}N_{\overline\ell\overline i}^{j}\overline\theta^\ell
 +\sum_{i,k}N_{k\overline i}^{j}\theta^k.
\]

下面小林昭七的定理阐明殆凯勒流形 $(M,J,g)$ 的 $T^{(1,0)}M$ 上典范联络与 $(M,g)$ 上列维--奇维塔联络之间的关系。

\begin{theorem}[小林昭七]\label{thm:canonical-equals-lichnerowitz}
\stfhIndex{I212}{nai^en^hei^si^wu^er^fu}{奈恩黑斯--伍尔夫（Nijenhuis–Wolf）}
假设 $(M,J,g)$ 是殆凯勒流形。则上一小节中典范联络在 $T^{(1,0)}M$ 上诱导的联络与利希内罗维茨联络重合，即 $\omega_i^j=\lambda_i^j$。特别地，$\Theta=\Theta'$，从而
\[
 N_{k\overline i}^{j}=0,
 \qquad N_{\overline\ell\overline i}^{j}
 =N_{\overline i\overline\ell}^{j}.
\]
\end{theorem}

与 $N_{\overline\ell\overline i}^{j}$ 相伴的 $(1,2)$-张量正是殆复结构 $J$ 的复化奈恩黑斯挠率张量。若令
$N_{\overline j\overline\ell\overline i}=N_{\overline\ell\overline i}^{j}$，则容易验证\eqref{eq:torsion-cyclic-symmetry}蕴含如下对称性质：
\begin{equation}
 N_{\overline j\overline\ell\overline i}
 =N_{\overline\ell\overline i\overline j}
 =N_{\overline i\overline j\overline\ell}.
 \label{eq:nijenhuis-cyclic-symmetry}
\end{equation}
把\eqref{eq:lichnerowitz-torsion}以及
\[
 \omega_i^j=\sum_\ell\omega_{i\ell}^j\theta^\ell
 +\sum_\ell\omega_{i\overline\ell}^j\overline\theta^\ell
\]
代入\eqref{eq:first-structure-unitary-connection}，并写成
\[
 d\theta^j=-\sum_{i,\ell}\omega_{i\ell}^j\theta^\ell\wedge\theta^i
 -\sum_{i,\ell}\omega_{i\overline\ell}^j\overline\theta^\ell\wedge\theta^i
 +\sum_{i,\ell}N_{\overline\ell\overline i}^{j}
 \overline\theta^\ell\wedge\overline\theta^i.
\]

最后考察殆凯勒流形上列维--奇维塔联络与典范联络的关系。设 $\nabla^{LC}$ 是列维--奇维塔联络。考虑乘法 $J:TM\to TM$ 的标准平均联络 $\nabla^{av}$：
\begin{equation}
 \nabla_X^{av}Y:=
 \frac{\nabla_X^{LC}Y+J^{-1}\nabla_X^{LC}(JY)}2
 =\nabla_X^{LC}Y-\frac12J(\nabla_X^{LC}J)Y.
 \label{eq:averaged-connection}
\end{equation}
于是有下面这个命题，它断言该联络成为典范联络。其证明可见 (KN96, Theorem 3.4)，或见 (Ga97) 第 2 节，特别是 (Ga97) 的 (2.2.10)。

\begin{theorem}\label{thm:averaged-connection-canonical}
假设 $(M,g,J)$ 是殆凯勒流形，即二形式 $\omega=g(J\cdot,\cdot)$ 是闭的。则平均联络 $\nabla^{av}$ 定义 $(M,g,J)$ 的典范联络；也就是说，该联络是 $J$-线性的、保持度量，并且其复化挠率具有型 $(0,2)$。
\end{theorem}

\begin{proof}
利用定义\eqref{eq:averaged-connection}，计算 $\nabla^{av}$ 的挠率 $T$，得到
\begin{equation}
\begin{split}
 T(X,Y)
 &=\nabla_X^{av}Y-\nabla_Y^{av}X-[X,Y]\\
 &=T^{LC}(X,Y)-\frac12J(\nabla_X^{LC}J)Y
 +\frac12J(\nabla_Y^{LC}J)X\\
 &=-\frac12J(\nabla_X^{LC}J)Y
 +\frac12J(\nabla_Y^{LC}J)X,
\end{split}
 \label{eq:averaged-connection-torsion}
\end{equation}
其中第二个等号使用了
$T^{LC}(X,Y)=\nabla_X^{LC}Y-\nabla_Y^{LC}X-[X,Y]=0$。另一方面，$d\omega=0$ 蕴含
\[
 \nabla_{JX}^{LC}J+J(\nabla_X^{LC}J)=0.
\]
（见 (Ga97) 的 (2.2.10)。）在\eqref{eq:averaged-connection-torsion}中取 $Y=JX$ 并代入这个等式，得到 $T(JX,X)=0$，这正是典范联络的定义性质。于是命题 7.1.11 蕴含挠率具有型 $(0,2)$。
\end{proof}

\subsection{辛流形 \texorpdfstring{$(M,\omega)$}{(M, omega)} 上的情形}
\label{subsec:canonical-connection-symplectic-manifolds}

现在考虑辛流形 $(M,\omega)$。

\begin{definition}\label{def:compatible-almost-complex-structure}
若殆复结构 $J$ 满足：
\begin{enumerate}[label=(\roman*)]
 \item （$J$-正）$\omega(u,Ju)\geq0$，且等号成立当且仅当 $u=0$；
 \item （$J$-埃尔米特）$\omega(Ju,Jv)=\omega(u,v)$，
\end{enumerate}
则称 $J$ 与 $\omega$ \emph{相容}。以 $\mathcal J_\omega$ 表示所有这种殆复结构的集合。
\end{definition}

极分解保证在任意辛流形 $(M,\omega)$ 上总存在一个相容殆复结构。

\begin{exercise}\label{ex:compatible-almost-complex-exists}
证明这个断言。
\end{exercise}

事实上，有格罗莫夫 (Gr85) 的下述一般引理。

\begin{lemma}[格罗莫夫]\label{lem:compatible-almost-complex-contractible}
集合 $\mathcal J_\omega$ 是一个非空、可缩、无穷维的弗雷歇流形。
\end{lemma}

对任意与 $\omega$ 相容的 $J$，双线性形式
\[
 g=\omega(\cdot,J\cdot)
\]
定义关于 $J$ 的埃尔米特度量，因而三元组 $(M,J,g)$ 按上述定义给出一个殆凯勒流形。依定义，这个度量 $g$ 的基本二形式成为 $\Phi=\omega$。为了突出辛形式 $\omega$，我们把这个殆凯勒流形记为 $(M,\omega,J)$。

从现在起，除非另有说明，在计算中我们主要使用定理 7.1.5 给出的典范殆埃尔米特联络，而不用殆凯勒流形 $(M,\omega,J)$ 的 $TM$ 上列维--奇维塔联络。不过我们指出，由定理 7.1.5，两个联络在全纯切丛 $T^{(1,0)}M$ 上诱导同一个相伴典范联络。

\begin{example}\label{ex:standard-euclidean-almost-kahler}
考虑标准欧氏辛向量空间
\[
 \mathbb R^{2n},\qquad
 \omega_0=\sum_{j=1}^n dq^j\wedge dp_j
\]
以及同构
\[
 \mathbb R^{2n}\longrightarrow\mathbb C^n;
 \quad(q^1,\ldots,q^n,p_1,\ldots,p_n)
 \longmapsto(z_1,\ldots,z_n),
 \qquad z_j=q^j+\sqrt{-1}p_j.
\]
由此，关于 $\mathbb C^n$ 上乘以 $i=\sqrt{-1}$ 所诱导的 $J_0$，欧氏内积 $g$ 写成 $g=\omega_0(\cdot,J_0\cdot)$；也就是说，$J_0$ 是由
\[
 J_0\left(\frac\partial{\partial q^j}\right)
 =\frac\partial{\partial p_j},\qquad
 J_0\left(\frac\partial{\partial p_j}\right)
 =-\frac\partial{\partial q^j}
\]
确定的自同态，其中 $j=1,\ldots,n$。注意，$\mathbb R^{2n}\cong\mathbb C^n$ 带有三个典范双线性形式：辛形式 $\omega_0$、\emph{欧氏内积} $g$ 和埃尔米特内积 $\langle\cdot,\cdot\rangle$。我们对这三者之间关系采用的约定是
\[
 \langle\cdot,\cdot\rangle
 =g(\cdot,\cdot)-i\omega_0(\cdot,\cdot).
\]
换言之，埃尔米特内积关于第一个变元复线性，关于第二个变元复反线性。
\end{example}
